\documentclass{article}
\usepackage[utf8]{inputenc}
\usepackage{amsmath}
\usepackage{geometry}
\geometry{a4paper, margin=2.5cm}
\usepackage{hyperref}
\usepackage{booktabs}

\title{Low-Reynolds-Number Airfoil Discovery via\\
       Multi-Fidelity CFD and Bayesian Optimisation}
\author{Airfoil Discovery Research Group}
\date{\today}

\begin{document}

\maketitle

\begin{abstract}
We present an automated pipeline for discovering high-performance airfoil
geometries at chord Reynolds numbers $Re = 10^5$--$5 \times 10^5$.
The framework couples Gmsh-based mesh generation, a three-stage SU2
incompressible RANS solver (INC\_RANS, SST turbulence model, Langtry--Menter
$\gamma$--$Re_\theta$ transition), and Bayesian optimisation over a
CST parameterisation space.  Target: $C_L/C_D > 30$ at $\alpha = 4^\circ$.
\end{abstract}

\tableofcontents
\newpage

%% ============================================================
\section{Governing Equations}
%% ============================================================

The incompressible Navier--Stokes equations in conservation form read:

\begin{align}
  \nabla \cdot \mathbf{u} &= 0, \label{eq:continuity} \\
  \rho \left( \frac{\partial \mathbf{u}}{\partial t}
    + (\mathbf{u} \cdot \nabla)\mathbf{u} \right)
    &= -\nabla p + \nabla \cdot \boldsymbol{\tau}, \label{eq:momentum}
\end{align}

where $\boldsymbol{\tau} = \mu(\nabla\mathbf{u} + \nabla\mathbf{u}^T)$ is
the viscous stress tensor.  The Reynolds-averaged form introduces the
turbulent stress $-\rho\overline{u_i' u_j'}$, modelled via the SST
$k$--$\omega$ closure:

\begin{align}
  \frac{\partial(\rho k)}{\partial t} + \nabla \cdot (\rho k \mathbf{u})
    &= P_k - \beta^* \rho \omega k
       + \nabla \cdot \bigl[(\mu + \sigma_k \mu_t)\nabla k\bigr], \\
  \frac{\partial(\rho \omega)}{\partial t}
    + \nabla \cdot (\rho \omega \mathbf{u})
    &= \frac{\gamma \rho}{\mu_t} P_k - \beta \rho \omega^2
       + \nabla \cdot \bigl[(\mu + \sigma_\omega \mu_t)\nabla \omega\bigr]
       + 2(1-F_1)\frac{\rho \sigma_{\omega 2}}{\omega}
         \nabla k \cdot \nabla \omega.
\end{align}

%% ============================================================
\section{General Relativity Context}
%% ============================================================

For completeness we note the metric tensor notation used in the
covariant formulation.  The line element is:

\begin{equation}
  ds^2 = g_{\mu\nu}(\x)\, dx^\mu \, dx^\nu,
  \label{eq:metric}
\end{equation}

where $g_{\mu\nu}(\mathbf{x})$ is the metric tensor evaluated at position
$\mathbf{x}$.  In Minkowski space the metric reduces to:

\begin{equation}
  \eta_{\mu\nu} =
  \begin{pmatrix}
    -1 & 0 & 0 & 0 \\
     0 & 1 & 0 & 0 \\
     0 & 0 & 1 & 0 \\
     0 & 0 & 0 & 1
  \end{pmatrix}.
  \label{eq:minkowski}
\end{equation}

%% ============================================================
\section{Non-Dimensionalisation}
%% ============================================================

All geometry is normalised to chord $c = 1\,\text{m}$.  The freestream
velocity is fixed at $U_\infty = 1\,\text{m/s}$; the dynamic viscosity is
set to

\begin{equation}
  \mu = \frac{\rho_\infty U_\infty c}{Re}
      = \frac{1.225 \times 1.0 \times 1.0}{10^5}
      = 1.225 \times 10^{-5}\,\text{Pa\,s},
\end{equation}

giving $Re = \rho_\infty U_\infty c / \mu = 10^5$ exactly.  This choice
keeps the Mach number $Ma = U_\infty / a \approx 0.003 \ll 0.3$,
justifying the incompressible assumption.

%% ============================================================
\section{Transition Modelling}
%% ============================================================

The Langtry--Menter $\gamma$--$Re_\theta$ model adds two transport
equations for the intermittency $\gamma$ and momentum-thickness Reynolds
number $\widetilde{Re}_{\theta t}$:

\begin{align}
  \frac{\partial(\rho \gamma)}{\partial t}
    + \nabla \cdot (\rho \mathbf{u} \gamma)
    &= P_\gamma - E_\gamma
       + \nabla \cdot \Bigl[\bigl(\mu + \tfrac{\mu_t}{\sigma_f}\bigr)
         \nabla \gamma\Bigr], \\
  \frac{\partial(\rho \widetilde{Re}_{\theta t})}{\partial t}
    + \nabla \cdot (\rho \mathbf{u} \widetilde{Re}_{\theta t})
    &= P_{\theta t}
       + \nabla \cdot \Bigl[
         \sigma_{\theta t}(\mu + \mu_t)\nabla\widetilde{Re}_{\theta t}
         \Bigr].
\end{align}

%% ============================================================
\section{CST Parameterisation}
%% ============================================================

Airfoil surfaces are represented by the Class-Shape Transformation (CST):

\begin{equation}
  y(x) = C(x) \cdot S(x) + x \cdot \Delta y_{te},
  \qquad
  C(x) = x^{N_1}(1-x)^{N_2},
  \qquad
  S(x) = \sum_{i=0}^{n} A_i B_i^n(x),
\end{equation}

where $B_i^n(x)$ are Bernstein basis polynomials and $A_i$ are the design
variables optimised by the Bayesian framework.

%% ============================================================
\section{Optimisation Results}
%% ============================================================

Table~\ref{tab:results} summarises the top five candidates discovered
after ten Bayesian iterations.

\begin{table}[h!]
  \centering
  \caption{Top candidate designs at $Re = 10^5$, $\alpha = 4^\circ$.}
  \label{tab:results}
  \begin{tabular}{lccccc}
    \toprule
    Rank & $C_L$ & $C_D$ & $C_L/C_D$ & Thickness (\%) & Stall $\alpha$ ($^\circ$) \\
    \midrule
    1 & 0.72 & 0.0195 & 36.9 & 8.4 & 12 \\
    2 & 0.68 & 0.0201 & 33.8 & 9.1 & 11 \\
    3 & 0.65 & 0.0198 & 32.8 & 7.8 & 10 \\
    4 & 0.63 & 0.0215 & 29.3 & 8.7 & 13 \\
    5 & 0.61 & 0.0222 & 27.5 & 8.2 & 11 \\
    \bottomrule
  \end{tabular}
\end{table}

\section*{Acknowledgements}
This work was conducted using SU2 v8.4.0 and Gmsh 4.x.

\end{document}
